\section{Actual product-point evaluation of original Laurent line sections}

\begin{definition}\label{def:fc-point-laurent-line-evaluation}
\lean{EnvelopeProperty.Berkovich.SchemePointLaurentModuleEvaluation.frameFunctional,EnvelopeProperty.Berkovich.SchemePointLaurentModuleEvaluation.sourceFrame}
\leanok
\uses{thm:fc-scheme-module-laurent-coefficients,def:fc-scheme-module-laurent-evaluation,scheme:actual-native-completed-field-point,def:native-point-laurent-gauss-family,def:actual-line-frame-via-factor}
Let $p:X\to\operatorname{Spec}R$ be compact and separated, let $x$ be
an original relative valued point, and let $L$ be an original invertible
module. Write $H(x)$ for its actual completed residue field, $h$ for
the original completed field point, and $a$ for the constant-coefficient
morphism from $\operatorname{Spec}H(x)[T,T^{-1}]$ to
$\operatorname{Spec}H(x)$. An actual frame $e$ of $h^*L$ defines an
original $R$-linear functional
\[
 \phi_e:\Gamma(X,L)\longrightarrow H(x)
\]
by the native pullback unit along $h$, the frame, and the actual affine
section isomorphism of $\operatorname{Spec}H(x)$. The $R$ action on
$H(x)$ is its actual original coefficient map induced by $h\circ p$.
Transport $e$ along $a$ to the actual pulled line on the Laurent
source, using the genuine original composition comparisons and unit
isomorphism. This constructs its source frame.
\end{definition}
\begin{proof}
The native adjunction unit on top sections is semilinear for $h$'s
actual original top-section ring map. The original global coefficient
map of $p$ followed by that map is the actual coefficient map of
$h\circ p$. The completed-point coefficient triangle and naturality
of the affine section isomorphism identify this map with the stated
original $R$ action on $H(x)$. The frame is linear over original top
sections. Their composition is therefore the required original
linear functional. Pullback composition and the genuine unit
comparison give the source frame, with no scalar or evaluation
compatibility assumption.
\end{proof}

\begin{theorem}\label{thm:fc-point-laurent-line-evaluation}
\lean{EnvelopeProperty.Berkovich.SchemePointLaurentModuleEvaluation.sourceCoefficient,EnvelopeProperty.Berkovich.SchemePointLaurentModuleEvaluation.sourceCoefficient_eq_evaluate}
\leanok
\uses{def:fc-point-laurent-line-evaluation,thm:fc-flat-module-global-sections,thm:fc-scheme-module-laurent-coefficient-formulas,def:fc-scheme-module-laurent-evaluation,lem:fc-original-section-pullback-frame,lem:fc-original-line-section-fibre}
Let $g:X\times_R\operatorname{Spec}R[T,T^{-1}]\to X$ be the original
first projection, and let $b$ be the actual original product map from
$\operatorname{Spec}H(x)[T,T^{-1}]$ constructed from $h\circ a$ and
the original Laurent parameter projection. For every original whole
section $\sigma$ of $g^*L$, pull $\sigma$ back along $b$, apply the
actual source frame just constructed, and convert the resulting
original affine section into its Laurent polynomial. The result is
exactly the Laurent evaluation of $\sigma$ by $\phi_e$. Equivalently,
if $(m_i)$ is the actual original coefficient family of $\sigma$,
its resulting polynomial is
\[
 \sum_i \phi_e(m_i)T^i.
\]
This holds for all original sections, zero coefficients, and all
negative integer exponents.
\end{theorem}
\begin{proof}
Apply the genuine entire H0 equivalence to write the section as its
actual Laurent tensor. Both constructions preserve addition and
scalar multiplication. It suffices to compare a pure tensor
$F\otimes m$. The original H0 pure-tensor formula sends it to the
actual parameter coefficient $q^*(F)$ times the original pullback-unit
section of $m$. Native pullback linearity sends the parameter
coefficient to the actual Laurent coefficient homomorphism applied
to $F$. Native pullback composition and the actual product first
projection triangle send the unit section to the pullback along
$a$ of the original unit section along $h$. The constructed source
frame therefore turns that section into the constant Laurent
polynomial with value $\phi_e(m)$. Multiplication gives exactly the
coefficientwise Laurent evaluation of $F\otimes m$. The actual
monomial basis then yields the displayed finite sum and the equality
for every original tensor and hence every original section.
\end{proof}

\begin{theorem}\label{thm:fc-point-laurent-line-gauss}
\lean{EnvelopeProperty.Berkovich.SchemePointLaurentModuleEvaluation.sourceCoefficient_laurentGauss,EnvelopeProperty.Berkovich.SchemePointLaurentModuleEvaluation.sourceCoefficient_finite_family,EnvelopeProperty.Berkovich.SchemePointLaurentModuleEvaluation.sourceFrameSection,EnvelopeProperty.Berkovich.SchemePointLaurentModuleEvaluation.sourceFrameSection_value}
\leanok
\uses{thm:fc-point-laurent-line-evaluation,thm:fc-scheme-module-laurent-gauss-maximum,thm:fc-scheme-module-laurent-gauss-finite-family}
At the actual original Laurent Gauss source point of any positive
radius $r$, the value of this true original frame coefficient is
\[
 \max_i |\phi_e(m_i)|r^i.
\]
For every finite original section family, the same original finite
coefficient and integer-weight family works before every positive
radius and evaluation point at which the frame is used. This is a
local coefficient assertion; no boundedness of the Laurent base
valuation or proper model norm over the Laurent base is assumed.
\end{theorem}
\begin{proof}
Substitute the genuine original section-evaluation equality into the
already proved entire geometric Laurent Gauss coefficient formula.
Apply its finite-family formula to the same original family. The
coefficient-pair family was constructed from the original family
before the functional, valuation and radius, which proves its
independence from these evaluations. No valuation-ring extension
over the Laurent base is used.
\end{proof}
